%% ****** Start of file aiptemplate.tex ****** %
%%
%%   This file is part of the files in the distribution of AIP substyles for REVTeX4.
%%   Version 4.1 of 9 October 2009.
%%
%
% This is a template for producing documents for use with
% the REVTEX 4.1 document class and the AIP substyles.
%
% Adapted from the AIP/REVTeX4-1 template to mirror the structure of
% jctc-draft.tex (ACS submission of the same manuscript). Body sections,
% appendix and macro definitions are kept in separate subfiles so that each
% can also be compiled/spell-checked on its own via the subfiles package.
%
% quantum.bib and paper2025_figs/ are shared as-is with jctc-draft.tex
% (referenced directly, not copied).

\documentclass[aip,jcp,graphicx,draft]{revtex4-1}

%%%%%%%%%%%%%%%%%%%%%%%%%%%%%%%%%%%%%%%%%%%%%%%%%%%%%%%%%%%%%%%%%%%%%
%% Font setup
%%%%%%%%%%%%%%%%%%%%%%%%%%%%%%%%%%%%%%%%%%%%%%%%%%%%%%%%%%%%%%%%%%%%%
\usepackage[T1]{fontenc}
\usepackage{lmodern}

%%%%%%%%%%%%%%%%%%%%%%%%%%%%%%%%%%%%%%%%%%%%%%%%%%%%%%%%%%%%%%%%%%%%%
%% Math support
%% (revtex4-1 only loads amsmath/amssymb if given as class options, so
%% pull them in explicitly here instead)
%%%%%%%%%%%%%%%%%%%%%%%%%%%%%%%%%%%%%%%%%%%%%%%%%%%%%%%%%%%%%%%%%%%%%
\usepackage{amsmath, amssymb}
\usepackage{braket}
\usepackage{bm}

%%%%%%%%%%%%%%%%%%%%%%%%%%%%%%%%%%%%%%%%%%%%%%%%%%%%%%%%%%%%%%%%%%%%%
%% Graphic inclusion
%% (the "graphicx" class option above only tells revtex4-1 which driver
%% to assume; it does not actually load the graphicx package.
%% [final] is needed because the "draft" class option is otherwise
%% forwarded to graphicx too, which would replace every figure with an
%% empty placeholder box instead of actually showing the image.)
%%%%%%%%%%%%%%%%%%%%%%%%%%%%%%%%%%%%%%%%%%%%%%%%%%%%%%%%%%%%%%%%%%%%%
\usepackage[final]{graphicx}

%%%%%%%%%%%%%%%%%%%%%%%%%%%%%%%%%%%%%%%%%%%%%%%%%%%%%%%%%%%%%%%%%%%%%
%% Split manuscript into subfiles, same layout as jctc-draft.tex
%%%%%%%%%%%%%%%%%%%%%%%%%%%%%%%%%%%%%%%%%%%%%%%%%%%%%%%%%%%%%%%%%%%%%
\usepackage{subfiles}

%% custom macros defined for this manuscript.
\usepackage{aip-draft_commands}

\begin{document}

% Use the \preprint command to place your local institutional report number
% on the title page in preprint mode.
%\preprint{}

\title{Time Evolution Simulation of a Hydrogen Atom Using Full Quantum Circuits}

\author{Hideo Takahashi}
\email{hideo.takahashi.qa@hitachi.com}
\affiliation{Managed \& Platform Services Business Division, Hitachi, Ltd.,
  292, Yoshida, Totsuka-ku, Yokohama, 244-0817, Japan}
\affiliation{Institute of Industrial Science, The University of Tokyo, 4-6-1, Komaba,
  Meguro-ku, Tokyo, 153-8505, Japan}

\author{Tatsuya Tomaru}
\affiliation{Next Research, Research \& Development Group, Hitachi, Ltd.,
  Kokubunji, Tokyo, 185-8601, Japan}
\affiliation{Institute of Industrial Science, The University of Tokyo, 4-6-1, Komaba,
  Meguro-ku, Tokyo, 153-8505, Japan}

\author{Toshiyuki Hirano}
\affiliation{Institute of Industrial Science, The University of Tokyo, 4-6-1, Komaba,
  Meguro-ku, Tokyo, 153-8505, Japan}
\affiliation{Tokyo Metropolitan College of Industrial Technology, 8-17-1, Minami-Senju, Arakawa-ku,
  Tokyo, 116-8523, Japan}

\author{Saisei Tahara}
\affiliation{Institute of Industrial Science, The University of Tokyo, 4-6-1, Komaba,
  Meguro-ku, Tokyo, 153-8505, Japan}

\author{Fumitoshi Sato}
\affiliation{Institute of Industrial Science, The University of Tokyo, 4-6-1, Komaba,
  Meguro-ku, Tokyo, 153-8505, Japan}

\date{\today}

\begin{abstract}

A quantum circuit for first-quantization-based Hamiltonian simulation, capable
of operating on a small-scale quantum computer emulator with fewer than 30
qubits was constructed. (1) For small model sizes, Uniformly Controlled Rotation
(UCR) circuits were found to be a preferable choice for calculating the Coulomb
potential compared with conventional arithmetic gates in terms of qubit count
and circuit depth. (2) As a guideline for obtaining optimal accuracy within the
given resources, we established a method for handling the singularity (pole) of
Coulomb potential and optimizing the discretization parameters. Regarding the
singularity of Coulomb potential, we demonstrated that divergence could be
avoided and errors minimized by determining an effective value of the potential
function from a simple geometric average of the potential within a minute
circular region centered on the pole. (3) We proposed a method to determine the
optimal spatial grid spacing for a given number of qubits for coordinates, by
evaluating the tradeoff between the spatial discretization error caused by grid
coarseness and the domain truncation error caused by a finite domain range. For
the time evolution step size, rather than relying on the formal Trotter error
bound that assumes the worst-case scenario, we observed that sampling-theorem
based bounds focusing on the maximum frequency component of the wave function
are sufficient to achieve optimal accuracy while avoiding an unnecessary
increase in the quantum circuit depth.

\end{abstract}

\maketitle %\maketitle must follow title, authors, abstract and \pacs

%%%%%%%%%%%%%%%%%%%%%%%%%%%%%%%%%%%%%%%%%%%%%%%%%%%%%%%%%%%%%%%%%%%%%
%% Start the main part of the manuscript here.
%%%%%%%%%%%%%%%%%%%%%%%%%%%%%%%%%%%%%%%%%%%%%%%%%%%%%%%%%%%%%%%%%%%%%
\subfile{aip-draft_1_intro.tex}
\subfile{aip-draft_2_method.tex}
\subfile{aip-draft_34_results_discussions.tex}

%%%%%%%%%%%%%%%%%%%%%%%%%%%%%%%%%%%%%%%%%%%%%%%%%%%%%%%%%%%%%%%%%%%%%
%% The conclusion section
%%%%%%%%%%%%%%%%%%%%%%%%%%%%%%%%%%%%%%%%%%%%%%%%%%%%%%%%%%%%%%%%%%%%%

\section{Conclusion}
In this study, we presented a circuit configuration capable of simulating the
full time evolution of a two-dimensional hydrogen atom model on a Qiskit
emulator. To the best of our knowledge, this is the first report to implement
grid-based time evolution calculations, including the Coulomb potential,
entirely as a quantum circuit and to confirm its operation on an emulator.

Regarding the handling of the Coulomb potential singularity, we proposed a
method using $\reff=\deltar/4$ as the effective distance to determine the
effective potential at the grid point containing the singularity in the 2D
model. Numerical experiments confirmed that this value minimizes the
energy error in the time-evolution simulation, and that it coincides with
the value derived from the geometric area average around the singularity,
clarifying the physical validity of this approach in a discretized space.

Furthermore, we established a procedure to determine the simulation parameters
that minimize the error for a given number of qubits ($\nb$). We showed that the
optimal spatial grid spacing, $\deltar$, can be determined from the trade-off
between the spatial discretization error and the domain truncation error.
Importantly, this spacing dictates the treatable range of
energy levels---from the ground state to higher-order excited states---for which
wave functions can be represented, thereby defining the effective dynamic range
of the simulation.

Moreover, regarding the time step, $\deltat$, of the Trotter decomposition, we
numerically confirmed that it can be determined not from the formal upper bound
of the Trotter error assuming the worst case, but from the requirement of the
sampling theorem (half-period condition) based on the maximum frequency
component of the wave function. Particularly when targeting states near the
ground state, we demonstrated that within the range satisfying this condition,
the spatial error becomes dominant over the temporal error (i.e., the error
reaches a lower bound), providing a guideline to avoid unnecessary increases in
quantum circuit depth.

These findings are expected to provide highly practical and powerful guidelines
for designing and executing quantum simulations using the grid method under
limited quantum resources.

\subfile{aip-draft_a_appendix.tex}

\begin{acknowledgments}
This research was partially supported by the UTokyo Quantum Initiative.
\end{acknowledgments}

\section*{Data and Software Availability Statement}

The data and source code underlying this study are openly available on GitHub at https://github.com/crsq-dev/report2026/releases.

% Create the reference section using BibTeX (apsrev4-1, bundled with revtex4-1):
\bibliographystyle{apsrev4-1}
\bibliography{quantum}

\end{document}
%
% ****** End of file aiptemplate.tex ******
